\chapter*{历史注记}
\phantomsection
\addcontentsline{toc}{chapter}{历史注记}

（注意 --- 罗马数字指本注记末尾的参考文献。）

随着 19 世纪“向量分析”的发展，积分向量值函数成为常见之事，但只要所涉及的只是取值于有限维空间的函数，这一运算并不构成问题。只是到了 Hilbert 的谱理论，才遇到自然导向更一般的积分概念的运算：这一理论实际上导致把 Hilbert 空间 H 上的每个连续 Hermite 型 \(\Phi(x,y)\) 对应于一个正交射影族 \((E(\lambda))_{\lambda\in\mathbf{R}}\) ，它具有这样的性质：对 H 的每一对向量 \((x,y)\) ，函数 \(\lambda\mapsto(E(\lambda)x|y)\) 是有界变差的，且 \(\Phi(x,y)=\int\lambda d(E(\lambda)x|y)\) ；若把 \(\Phi\) 与使 \(\Phi(x,y)=(Ax|y)\) 成立的 Hermite 算子 A 相联系，人们便很想把前一公式写成 \(A=\int\lambda dE(\lambda)\) 。但只是从大约 1935 年起，在 Bochner 引进了取值于 Banach 空间的函数的（“强”）积分之后，人们才开始致力于定义向量值函数的积分（或关于向量测度的积分），使之能够合法地写出诸如前一公式那样的公式。这一推广主要是由 Gelfand（III）、Dunford 与 Pettis（IV）和（V）完成的；他们的结果是对 Banach 空间叙述的，但可毫无困难地推广到更一般的局部凸空间。

把一个体积分解成“切片”、并把该体积上的积分化为每个切片上的积分、随后做一次单独积分的想法，在分析学中自无穷小演算的发端起即已被使用（Cavalieri 的“不可分元演算”不过是这一原理的最初轮廓，它甚至可以追溯到 Archimedes（参见第 IV 卷第 I--III 章的历史注记））。但在经典的应用中，“切片”总是具有非常特殊、非常正则的性质（最常见的是解析曲面的开子集，解析地依赖于一个参数）；在没有一般的积分理论的情况下，也很难有别的情形。测度分解的一般问题由 von Neumann 于 1932 年在遍历理论的关联中提出并解决（I）；大约在同一时间（并且独立地），Kolmogoroff 在为概率论奠定公理化基础时，被引向以一般的方式定义“条件概率”概念并证明其存在性，这个问题本质上等价于测度的分解问题（II）。

